% !TEX program = xelatex
% Chinese report layout: par=indent, secnum=zh, 13pt, kai, text=times,
% subtitle and info line in the title block.
\documentclass[lang=zh, thmnames=native, par=indent, secnum=zh, 13pt, kai,
  text=times]{akizumi-note}
\title{2013 年臺灣毒澱粉事件}
\author{林子祐}
\akizumisetup{course=食品安全, subtitle={食品安全風險溝通失敗之案例研究},
  info={學號=B10303083, 系級=資工四, 期末報告}}
\begin{document}
\maketitle

\section{緒論}\label{sec:intro}
食品安全涉及公眾健康與社會信任，其中風險溝通（Risk Communication）是連接科學評估、政策管理與消費者認知的關鍵環節。The first paragraph after a heading is indented too.

第二段同樣縮排兩個字，段落之間沒有額外間距。\textbf{粗體楷書}與\emph{強調}。

\subsection{研究背景}\label{sec:bg}
本研究見\cref{sec:intro}與\cref{sec:bg}，以及\cref{thm:a}和\cref{eq:a}。
\begin{equation}\label{eq:a} Y_i = f(A_i) + \varepsilon_i \end{equation}

\begin{theorem}\label{thm:a}
定理內文。
\end{theorem}

\subsubsection{細項}
\begin{itemize}
  \item 第一項
  \item 第二項
\end{itemize}

\setcounter{section}{22}
\section{結論}
\setcounter{subsection}{11}
\subsection{十二}

\ExplSyntaxOn
\str_if_eq:eeF { \thesection } { 二十三 } { \ERROR }
\str_if_eq:eeF { \thesubsection } { （十二） } { \ERROR }
\str_if_eq:eeF { \akizumizhnum{10}\akizumizhnum{11}\akizumizhnum{40}\akizumizhnum{105} } { 十十一四十105 } { \ERROR }
\ExplSyntaxOff
\ifdim\parindent=26pt \else\ERROR\fi  % 2em at 13pt
\end{document}
